\documentclass[UTF8,12pt]{ctexart}
\usepackage[a4paper,margin=24mm]{geometry}
\usepackage{amsmath,amssymb,bm,graphicx,adjustbox}
\newcommand{\fitmath}[1]{\begin{adjustbox}{max width=\linewidth}$\displaystyle #1$\end{adjustbox}}
\setlength{\parindent}{0pt}
\setlength{\parskip}{.7em}
\sloppy
\allowdisplaybreaks
\begin{document}
\section*{2016 年考研数学三 · 参考答案}
\subsection*{原 PDF 第 20 页}

\subsection*{2016年真题参考答案}

\subsection*{一、选择题}

（1）B。 （2）D。 （3）B。 （4）A。 （5）C。 （6）C。 （7）A。 （8）C。


\subsection*{二、填空题}

（9）\(\displaystyle 6\)。 （10）\(\displaystyle \sin1-\cos1\)。 （11）\(\displaystyle -\mathrm dx+2\mathrm dy\)。 （12）\(\displaystyle \dfrac{\displaystyle 1}{\displaystyle 3}-\dfrac{\displaystyle 2}{\displaystyle 3e}\)。 （13）\(\displaystyle \lambda^4+\lambda^3+2\lambda^2+3\lambda+4\)。 （14）\(\displaystyle \dfrac{\displaystyle 2}{\displaystyle 9}\)。


\subsection*{三、解答题}

（15）\(\displaystyle e^{1/3}\)。


（16）（I）\(\displaystyle Q(p)=1200-10p\)。（II）当 \(\displaystyle p=100\) 万元时，边际收益为80万元。其经济意义为：当 \(\displaystyle p=100\) 万元，\(\displaystyle Q=200\) 件时，销售第201件商品所得的收益为80万元。


（17）\(\displaystyle f^{\prime}(x)=\begin{cases}4x^2-2x,\allowbreak &0<x<1,\allowbreak \\2x,\allowbreak &x\ge1.\end{cases}\) 最小值 \(\displaystyle f(\dfrac{\displaystyle 1}{\displaystyle 2})=\dfrac{\displaystyle 1}{\displaystyle 4}\)。


（18）\(\displaystyle f(x)=-\dfrac{\displaystyle e^{-x}+e^x}{\displaystyle 2}\)。


（19）和函数 \(\displaystyle s(x)=\begin{cases}x\ln(\dfrac{\displaystyle 1+x}{\displaystyle 1-x})+\ln(1-x^2),\allowbreak &x\in(-1,\allowbreak 1),\allowbreak \\2\ln2,\allowbreak &x=\pm1.\end{cases}\) 收敛域为 \(\displaystyle [-1,\allowbreak 1]\)。


（20）（I）\(\displaystyle a=0\)。（II）通解为 \(\displaystyle k(0,\allowbreak 1,\allowbreak -1)^{\mathrm T}+(1,\allowbreak -2,\allowbreak 0)^{\mathrm T}\)，其中 \(\displaystyle k\) 为任意常数。


（21）（I）\(\displaystyle A^{99}=\begin{pmatrix}2^{99}-2&1-2^{99}&2-2^{98}\\2^{100}-2&1-2^{100}&2-2^{99}\\0&0&0\end{pmatrix}\)。（II）\(\displaystyle \beta_1=(2^{99}-2)\alpha_1+(2^{100}-2)\alpha_2\)，\(\displaystyle \beta_2=(1-2^{99})\alpha_1+(1-2^{100})\alpha_2\)，\(\displaystyle \beta_3=(2-2^{98})\alpha_1+(2-2^{99})\alpha_2\)。


（22）（I）\(\displaystyle f(x,\allowbreak y)=\begin{cases}3,\allowbreak &0<x<1,\allowbreak \ x^2<y<\sqrt x,\allowbreak \\0,\allowbreak &\text{其他}.\end{cases}\)（II）\(\displaystyle U\) 与 \(\displaystyle X\) 不相互独立，理由见解析。（III）\(\displaystyle F(z)=\begin{cases}0,\allowbreak &z\le0,\allowbreak \\\dfrac{\displaystyle 3}{\displaystyle 2}z^2-z^3,\allowbreak &0<z\le1,\allowbreak \\2(z-1)^{3/2}-\dfrac{\displaystyle 3}{\displaystyle 2}z^2+3z-1,\allowbreak &1<z\le2,\allowbreak \\1,\allowbreak &z>2.\end{cases}\)


（23）（I）\(\displaystyle f_T(t)=\begin{cases}\dfrac{\displaystyle 9t^8}{\displaystyle \theta^9},\allowbreak &0<t<\theta,\allowbreak \\0,\allowbreak &\text{其他}.\end{cases}\)（II）当 \(\displaystyle a=\dfrac{\displaystyle 10}{\displaystyle 9}\) 时，\(\displaystyle E(aT)=\theta\)。


\clearpage
\end{document}
